\subsection{端点で大きくなる関数の積分}

指数を正の実数へ広げるため，$p,q>0$として
\[
 x^{p-1}(1-x)^{q-1}
\]
を考えたい。例えば$p=\frac12$なら$x^{p-1}=\frac1{\sqrt x}$となり，$x=0$で値が定まらない。閉区間で連続な関数の積分という，これまでの定義をそのまま使うことはできない。

そこで，問題のある端点をいったん避けて積分し，その端点へ近づける。$0<c<1$を一つ固定して
\begin{equation}\label{式:両端広義積分の定義}
 \int_0^1 f(x)\dd x
 =\lim_{\varepsilon\downarrow0}\int_{\varepsilon}^{c}f(x)\dd x
  +\lim_{\delta\downarrow0}\int_c^{1-\delta}f(x)\dd x
\end{equation}
と定める。ただし，右辺の\emph{二つの極限がそれぞれ有限の値を持つ場合}に限る。このような積分を\textbf{広義積分}といい，有限の値を持つことを収束するという。$c$を別の内点に変えても，その二点間の通常の積分が一方から他方へ移るだけなので，和の値は変わらない。

両端へ同じ速さで近づける$\int_{\varepsilon}^{1-\varepsilon}$の極限だけでは，この定義を満たしたことにはならない。一方の端で正に発散し，他方の端で負に発散するものが打ち消し合うおそれがあるためである。この章のベータ積分は非負関数だが，後にも使えるよう，両端を独立に調べる定義にしておく。

まず，$p>0$なら
\[
 \int_{\varepsilon}^{c}x^{p-1}\dd x
 =\frac{c^p-\varepsilon^p}{p}
 \longrightarrow\frac{c^p}{p}
\]
である。$p$が$1$より小さくても，$\varepsilon^p\to0$となるので収束する。関数の高さが無限に大きくなることと，面積が無限に大きくなることは同じではない。

ベータ積分の左半分では，$x\in\left[0,\frac12\right]$において$(1-x)^{q-1}$が連続で，ある正の定数$M$以下である。従って
\[
 0\leqq x^{p-1}(1-x)^{q-1}\leqq Mx^{p-1}.
\]
下端$\varepsilon$を小さくすると積分値は増加するが，上の評価により$\frac{M}{p}\left(\frac12\right)^p$以下にとどまる。単調で上に有界な数列の極限と同じ考え方で，この極限は有限に存在する。右半分も$u=1-x$と置けば$q>0$によって同様に収束する。したがって$p,q>0$なら，両端の広義積分がともに定まる。

\begin{definition}{ベータ関数}
$p,q>0$に対して
\begin{equation}\label{式:ベータ関数の定義}
 B(p,q)=\int_0^1x^{p-1}(1-x)^{q-1}\dd x
\end{equation}
と定める。端点で関数が定まらない場合には，式\eqref{式:両端広義積分の定義}による広義積分とする。
\end{definition}

整数の場合との対応は，指数に$1$のずれがある。式\eqref{式:整数ベータ積分}は
\[
 B(m+1,n+1)=\frac{m!n!}{(m+n+1)!}
\]
と書ける。$B(m,n)$の指数は$m-1,n-1$であり，$m,n$そのものではない。記号の引数と被積分関数の指数を取り違えないようにする。

